% §6 Open Problems
The two-floor model leaves deliberately open problems that we list with their intended next
steps.

\paragraph{Standardised cessation and perceived-V2V measurement.}
Industry measures barge-in timing ad hoc; there is no shared definition and no protocol for the
perceived variant (audio actually reaching the speaker, including transport and devices).
A standardised protocol for \emph{perceived} response onset and cessation -- analogous to the
reporting standards that exist for static benchmarks -- would let the floors of
Theorem~\ref{thm:cessation_floor} be compared across systems and labs.
Our smoke anchor is an approximation; the protocol is the full measurement campaign.

\paragraph{Guardrail re-arming under reentry.}
An execution-persistent orchestrator re-enters the model on every tool completion.
Each re-entry is a fresh generation window and therefore a fresh place a guardrail must arm
\citep{dong2024building}.
The safety-deadline question -- where a guardrail can sit inside the floor without inflating it
-- becomes a per-reentry schedulability problem with the same shared-dependency structure as
Proposition~\ref{prop:floor_dependency}.
This connects the present paper to the open streaming-safe guardrail gap and its
linear-cost certification question.

\paragraph{Forgetting and context compaction across long sessions.}
Long stateful sessions grow the prefill contribution to the onset floor over time.
Managing what the agent retains -- compaction, eviction, external memory -- is a
\emph{deadline-aware} policy in the same sense as buffer depth is here
\citep{xiao2023streamingllm,packer2023memgpt}.
A ``forgetting floor'' that ties retained context to the onset budget is the natural companion
result.

\paragraph{Cancellable vs irrevo\-cable tools.}
The ``tool persistence under cancellation is safe'' headline of
Section~\ref{sec:analysis} is \emph{protocol}-safe: it guarantees a provider-valid history
(Lemma~\ref{lem:serializability}).
It is not automatically \emph{business}-safe.
For irrevocable, side-effecting tools -- payment, transfer, authorization -- a resumed context
that re-runs the same tool risks double-execution or double-charge.
Idempotency keys, revocation, and explicit re-authorisation are therefore prerequisites for
persisting business-critical tools across an interruption.
This is a separate concern from the type A/B/C classification: tool typing describes when
a result is needed for audio, not whether re-running the persisted call is safe.

\paragraph{Learned turn-end prediction.}
Reacting to an interruption is the worst case; predicting it is the improvement.
The companion paper's predictive-stop ceiling and this paper's reactive floor delimit the
intervention window, and learning to shrink it is an open empirical question.

\paragraph{Predictive dispatch across tool types.}
Reclassifying a tool from type A to type B at runtime -- dispatching it off a partial
transcript before the turn is final -- is the engine behind deployed pre-EOU dispatch and
prefetch systems \citep{specrag2024,agentsnotenough2024,lin2026fdbv3}.
A composition rule that states when such pre-dispatch is safe, and what it costs in accuracy,
remains open.
